\problemname{Binära palindrom}
\noindent

Mikael studerar intressanta positiva heltal, och vill nu undersöka tal som är ett \emph{binärt palindrom}. 
För att formellt beskriva detta definierar vi följande funktioner:

\[
\text{bin}(n) = \text{den binära representationen av heltalet } n \text{ utan inledande nollor},
\]

\[
\text{reverse}(s) = \text{strängen } s \text{ skriven baklänges}.
\]

Ett heltal $n$ är alltså ett \emph{binärt palindrom} om
\[
\text{bin}(n) = \text{reverse}(\text{bin}(n)).
\]

Notera att ett palindrom är en sträng som läses likadant framifrån som bakifrån.


Mikael gillar de binära palindromer som innehåller exakt $k$ ettor i dess binära representation, 
och vill nu veta hur många sådana tal som finns inom intervallet $[l,r]$.

Givet tre heltal $l$, $r$ och $k$, bestäm hur många heltal $x$ som uppfyller:
\begin{enumerate}
    \item Talet $x$ är mellan $l$ och $r$, alltså att $l \leq x \leq r$,
    \item Talet $x$ är ett tal som Mikael gillar (dvs.\ ett binärt palindrom),
    \item Binärrepresentationen av $x$ innehåller exakt $k$ ettor.
\end{enumerate}

% Uppgiften är alltså att räkna antalet tal i intervallet $[l, r]$ som är ett binärt palindrom och som dessutom har exakt $k$ ettor i sin binära form.


\section*{Indata}
Indatan består av en rad med tre heltal $l, r, k$ ($1 \leq l \leq r \leq 10^{18}, 1 \leq k < 60$).

\section*{Utdata}
Skriv ut ett heltal, antalet tal i intervallet $[l, r]$ som är binära palindromer och har exakt $k$ ettor i sin binära representation.


\section*{Poängsättning}
Din lösning kommer att testas på en mängd testfallsgrupper.
För att få poäng för en grupp så måste du klara alla testfall i gruppen.

\noindent
\begin{tabular}{| l | l | p{12cm} |}
  \hline
  \textbf{Grupp} & \textbf{Poäng} & \textbf{Gränser} \\ \hline
  $1$    & $20$       & $l,r \leq 10^9$ \\ \hline
  $2$    & $20$       & $k \leq 4$ \\ \hline
  $3$    & $60$       & Inga ytterligare begränsningar. \\ \hline
\end{tabular}

\section*{Förklaring av exempelfall 1}

Vi söker tal i intervallet $[7, 10]$ som både är binära palindromer och har exakt två ettor i sin binära representation.

Det finns endast 4 tal i intervallet:

\begin{align*}
7 &= 111_2 \quad (\text{palindrom, tre ettor}) \\
8 &= 1000_2 \quad (\text{ej palindrom}) \\
9 &= 1001_2 \quad (\text{palindrom, två ettor}) \\
10 &= 1010_2 \quad (\text{ej palindrom})
\end{align*}

Det enda talet som uppfyller båda villkoren är $9$, så svaret är $1$.

\section*{Förklaring av exempelfall 2}

Vi söker tal i intervallet $[1, 6]$ som både är binära palindromer och har exakt två ettor i sin binära representation.

Det finns endast 6 tal i intervallet:

\begin{align*}
1 &= 1_2 \quad (\text{palindrom, en etta}) \\
2 &= 10_2 \quad (\text{ej palindrom}) \\
3 &= 11_2 \quad (\text{palindrom, två ettor}) \\
4 &= 100_2 \quad (\text{ej palindrom}) \\
5 &= 101_2 \quad (\text{palindrom, två ettor}) \\
6 &= 110_2 \quad (\text{ej palindrom})
\end{align*}

Talen som uppfyller båda villkoren är $3$ och $5$, så svaret är $2$.
